\subsection*{This analysis}

Every pair was fitted under a single QC arm, \textbf{001}: lenient per-variant
filters (MAF $\geq 0.01$, imputation INFO $\geq 0.9$, $\chi^2 \leq 80$,
allele-frequency discordance against the reference $\leq 0.20$), long-range LD
regions and the MHC \emph{retained}, and the LD-consistency screen \emph{on},
four rounds. That arm was chosen because the companion \texttt{ldl-cad} study
found the screen to be the load-bearing switch of the three; fixing the other
two isolates it.

The screen is run once per trait on the lenient variant set, and the resulting
masks are intersected across the pair --- not re-run per pair. Re-screening per
pair would let the screen see a different variant set each time, which entangles
what the per-variant filters did with what the screen was given.

Traits, source files, sample sizes, ancestries and known dataset defects are
registered in \texttt{../\_lib/traits.toml}. Sources that were tried and
rejected are recorded with their symptoms in \texttt{../\_lib/datasets.md}; read
that before adding a dataset, because two of the three defects recorded there
are invisible until a fit has already started.

This is one QC arm on one LD reference. It measures how the two summaries
behave across these pairs; it is not a comparison of QC strategies, and the
factorial that would support one is in \texttt{benchmarks/qc\_factorial.py}.
